\section{Distributional Semantics}
\label{sec:distributional-semantics}
\slidesrc{Distributional semantics: 1--27; Semantics (long version): 1--6}

With this section we move up from syntax to \textbf{semantics}: \emph{what does it
mean?} (Section~\ref{sec:three-levels}). The first and today most successful approach
does not define meanings at all: it derives them from \emph{how words are used} in large
corpora. This is the idea behind word embeddings and, ultimately, behind LLMs. The
slides are the first 13 slides of the \emph{Distributional Semantic Models} tutorial by
Stefan Evert, Alessandro Lenci, Marco Baroni and Gabriella Lapesa
(\url{http://wordspace.collocations.de}).

This section covers:
\begin{enumerate}[nosep]
  \item the distributional hypothesis (Section~\ref{sec:distributional-hypothesis});
  \item the co-occurrence matrix and a thought experiment that shows its power
        (Section~\ref{sec:cooccurrence-matrix});
  \item the geometric view: words as vectors, similarity as angle
        (Section~\ref{sec:geometric-interpretation}).
\end{enumerate}

\subsection{The Distributional Hypothesis}
\label{sec:distributional-hypothesis}
\slidesrc{Distributional semantics: 3--16; Semantics (long version): 3--5}

\subsubsection{Meaning and distribution}
\slidesrc{Distributional semantics: 4--7}

Four famous quotations, each followed by what it means for us:

\begin{itemize}
  \item ``\emph{Die Bedeutung eines Wortes liegt in seinem Gebrauch}'' (``the meaning of
        a word is its use'', Ludwig Wittgenstein).\\
        $\Rightarrow$ \textbf{meaning = use = distribution in language}.
  \item ``\emph{You shall know a word by the company it keeps!}'' (J.\,R.~Firth, 1957).\\
        $\Rightarrow$ \textbf{distribution = collocations} = habitual word combinations
        (Section~\ref{sec:corpora}).
  \item \textbf{Distributional hypothesis}: difference of meaning correlates with
        difference of distribution (Zellig Harris, 1954).\\
        $\Rightarrow$ we can measure a \textbf{semantic distance} between words.
  \item ``\emph{What people know when they say that they know a word is not how to recite
        its dictionary definition -- they know how to use it [\ldots] in everyday
        discourse}'' (George Miller, 1986).
\end{itemize}

\begin{nbdef}{Distributional hypothesis}
Words that occur in similar contexts tend to have similar meanings. Conversely, the more
the contexts of two words differ, the more their meanings differ.
\end{nbdef}

\subsubsection{What is the meaning of ``bardiwac''?}
\slidesrc{Distributional semantics: 8--16}

You have never seen the word \emph{bardiwac}. Read these sentences, one at a time:
\begin{itemize}[nosep]
  \item He handed her her glass of bardiwac.
  \item Beef dishes are made to complement the bardiwacs.
  \item Nigel staggered to his feet, face flushed from too much bardiwac.
  \item Malbec, one of the lesser-known bardiwac grapes, responds well to Australia's
        sunshine.
  \item I dined off bread and cheese and this excellent bardiwac.
  \item The drinks were delicious: blood-red bardiwac as well as light, sweet Rhenish.
\end{itemize}

\begin{nbintuition}[what you just did]
After a few sentences you know that \textbf{bardiwac is a heavy red alcoholic beverage
made from grapes}, i.e.\ a red wine (the word is invented; it stands for
\emph{claret}). Nobody gave you a definition: you inferred it from the
\emph{company the word keeps}: \emph{glass of}, \emph{too much}, \emph{grapes},
\emph{drinks}, \emph{blood-red}, \emph{beef dishes}. That is the distributional hypothesis
in action, and it is how we learn most of our vocabulary.
\end{nbintuition}

The examples are hand-picked and edited, of course. But in a corpus like the British
National Corpus (BNC) there is at least as much relevant information. The slides show a
\emph{word sketch} of the real word (BNC frequency 230): grouped by grammatical relation,
its most typical companions are verbs like \emph{uncork}, \emph{gulp}, \emph{sip},
\emph{drink}, \emph{swig}; coordinated nouns like \emph{burgundy}, \emph{Sancerre},
\emph{scampi}, \emph{ruby}; containers like \emph{bottle}, \emph{glass}, \emph{jug},
\emph{goblet}; modifiers like \emph{full-bodied}, \emph{blood-red}, \emph{ready-to-drink}.
A computer can collect these statistics automatically.

\subsection{The Co-occurrence Matrix}
\label{sec:cooccurrence-matrix}
\slidesrc{Distributional semantics: 17--21}

\subsubsection{A thought experiment: deciphering hieroglyphs}
\slidesrc{Distributional semantics: 17--20}

Imagine an unknown language. We know the meaning of some nouns (in brackets) but not of
\emph{dog}, and we do not know any of the verbs. For each noun we count how often it
appears as the \emph{object} of each verb in a large corpus:

\begin{center}
\small
\begin{tabular}{@{}l l rrrrrr@{}}
\toprule
 & & \textbf{get} & \textbf{sij} & \textbf{ius} & \textbf{hir} & \textbf{iit} & \textbf{kil} \\
\midrule
(knife)  & naif  & 51  & 20 & 84 & 0  & 3  & 0  \\
(cat)    & ket   & 52  & 58 & 4  & 4  & 6  & 26 \\
???      & dog   & 115 & 83 & 10 & 42 & 33 & 17 \\
(boat)   & beut  & 59  & 39 & 23 & 4  & 0  & 0  \\
(cup)    & kap   & 98  & 14 & 6  & 2  & 1  & 0  \\
(pig)    & pigij & 12  & 17 & 3  & 2  & 9  & 27 \\
(banana) & nana  & 11  & 2  & 2  & 0  & 18 & 0  \\
\bottomrule
\end{tabular}
\end{center}

Which known word is \emph{dog} most similar to? Comparing rows (the similarity measure is
explained in Section~\ref{sec:geometric-interpretation}):
\[
  \text{sim}(\text{dog}, \text{naif}) = 0.770, \qquad
  \text{sim}(\text{dog}, \text{pigij}) = 0.939, \qquad
  \text{sim}(\text{dog}, \text{ket}) = 0.961.
\]
\emph{Dog} behaves most like \emph{cat}, then \emph{pig}, and much less like \emph{knife}:
it is probably an animal, a pet. We concluded this without understanding a single verb.

\subsubsection{English as seen by the computer}
\slidesrc{Distributional semantics: 21}

The ``unknown language'' was English with scrambled spelling: the verbs are
\textbf{get, see, use, hear, eat, kill}, and the counts are real verb--object
co-occurrence counts from the British National Corpus. For a computer, every language
looks like hieroglyphs: it only has the counts.

\begin{nbdef}{Co-occurrence matrix and distributional semantic model}
A \textbf{co-occurrence matrix} $M$ has one row per \textbf{target word} and one column
per \textbf{context} (here: verbs taking the word as object; in general: nearby words,
syntactic relations, documents, \ldots). The cell $M_{w,c}$ counts how often $w$ occurs
with context $c$. The row $\mathbf{x}_w$ is the \textbf{distributional vector} of $w$.
A \textbf{distributional semantic model} (DSM) represents the meaning of each word by
such a vector.
\end{nbdef}

\begin{nbwarning}[how the similarity values were computed]
The slides do not say which similarity gives 0.770, 0.939 and 0.961. They are the
\textbf{cosine similarity} (Section~\ref{sec:geometric-interpretation}) of the rows after
replacing each count $f$ by $\log(1+f)$. On the raw counts the cosine gives different
values:
\begin{center}
\small
\begin{tabular}{@{}l cccccc@{}}
\toprule
similarity to \emph{dog} & knife & cat & boat & cup & pig & banana \\
\midrule
cosine, raw counts          & 0.552 & \textbf{0.923} & 0.915 & 0.832 & 0.676 & 0.629 \\
cosine, $\log(1+f)$         & 0.770 & \textbf{0.961} & 0.828 & 0.850 & 0.939 & 0.785 \\
\bottomrule
\end{tabular}
\end{center}
With raw counts \emph{boat} is almost as close to \emph{dog} as \emph{cat}, and
\emph{pig} is far: the large counts of the very frequent verb \emph{get} dominate
everything. The logarithm dampens large counts, so the rarer but more telling contexts
(\emph{hear}, \emph{eat}, \emph{kill}) matter too, and the two animals come out on top.
Real DSMs always apply such a weighting (log, TF-IDF-like association scores) before
comparing vectors.
\end{nbwarning}

\subsection{Geometric Interpretation}
\label{sec:geometric-interpretation}
\slidesrc{Distributional semantics: 22--27; Semantics (long version): 6}

\subsubsection{Words as points}
\slidesrc{Distributional semantics: 22--24}

The row vector $\mathbf{x}_{\text{dog}} = (115, 83, 10, 42, 33, 17)$ describes how
\emph{dog} is used in the corpus. It can be seen as the coordinates of a \textbf{point}
in an $n$-dimensional Euclidean space ($n = 6$ here, one dimension per context). We can
draw two dimensions, \emph{get} and \emph{use}: $\mathbf{x}_{\text{dog}} = (115, 10)$,
$\mathbf{x}_{\text{cat}} = (52, 4)$, $\mathbf{x}_{\text{knife}} = (51, 84)$,
$\mathbf{x}_{\text{boat}} = (59, 23)$.

\begin{center}
\begin{tikzpicture}
\begin{axis}[width=10cm, axis equal image, xmin=0, xmax=130, ymin=0, ymax=105,
    xlabel={get}, ylabel={use}, title={Two dimensions of the English V--Obj DSM},
    title style={font=\small}, tick label style={font=\scriptsize}, label style={font=\small},
    axis lines=left, clip=false]
  % unit-length arc (radius 80, as in the slides)
  \addplot[domain=0:90, samples=60, nbgray, thin] ({80*cos(x)}, {80*sin(x)});
  \draw[-{Stealth}, thick, nbgreen]  (axis cs:0,0) -- (axis cs:51,84);
  \draw[-{Stealth}, thick, nbblue]   (axis cs:0,0) -- (axis cs:59,23);
  \draw[-{Stealth}, thick, black]    (axis cs:0,0) -- (axis cs:52,4);
  \draw[-{Stealth}, thick, nbred]    (axis cs:0,0) -- (axis cs:115,10);
  \fill[nbgreen] (axis cs:51,84) circle (2pt) node[above left, font=\small] {knife};
  \fill[nbblue]  (axis cs:59,23) circle (2pt) node[above, font=\small] {boat};
  \fill[black]   (axis cs:52,4)  circle (2pt) node[above left, font=\small] {cat};
  \fill[nbred]   (axis cs:115,10) circle (2pt) node[above, font=\small] {dog};
  % normalised points on the arc
  \draw[nbgreen, thick] (axis cs:41.5,68.4) circle (2.5pt);
  \draw[nbblue, thick]  (axis cs:74.5,29.1) circle (2.5pt);
  \draw[nbred, thick]   (axis cs:79.7,7.0)  circle (2.5pt);
  % angle between knife and dog
  \addplot[domain=5:58.7, samples=40, dashed, thick] ({48*cos(x)}, {48*sin(x)});
  \node[font=\small] at (axis cs:48,50) {$\alpha$};
\end{axis}
\end{tikzpicture}
\captionof{figure}{Words as vectors in the (get, use) plane. Circles on the arc: the same
vectors normalised to the same length (only the direction is kept). \emph{Dog} and
\emph{cat} point in almost the same direction even though they are far apart as points.}
\end{center}

A first idea: \textbf{similarity = spatial proximity}, measured with the Euclidean
distance. But the location of a point depends on the \textbf{frequency} of the noun:
\emph{dog} is about 2.7 times more frequent than \emph{cat}, so all its counts are larger
and its point lies much farther from the origin.

\begin{nbexample}{Euclidean distance is misleading}
\begin{align*}
  d(\text{dog}, \text{cat})  &= \sqrt{(115-52)^2 + (10-4)^2}  = \sqrt{3969 + 36}  \approx 63.3\\
  d(\text{dog}, \text{boat}) &= \sqrt{(115-59)^2 + (10-23)^2} = \sqrt{3136 + 169} \approx 57.5
\end{align*}
By Euclidean distance, \emph{boat} is closer to \emph{dog} than \emph{cat} is: a
consequence of frequency, not of meaning.
\end{nbexample}

\subsubsection{Direction, not location}
\slidesrc{Distributional semantics: 25--27}

A vector can also be seen as an \textbf{arrow from the origin}. Its \emph{length} mostly
reflects how frequent the word is; its \emph{direction} reflects \emph{the proportions}
in which the word occurs with the different contexts, i.e.\ how it is used. So the
direction is more important than the location. Two equivalent solutions:
\begin{itemize}[nosep]
  \item use the \textbf{angle} $\alpha$ between two vectors as a distance measure;
  \item or \textbf{normalise} the length $\norm{\mathbf{x}}$ of every arrow to the same
        value, and then compare points (the circles on the arc in the figure).
\end{itemize}

\begin{nbdef}{Cosine similarity}
The \textbf{cosine similarity} of two vectors is the cosine of the angle between them:
\[
  \cos\alpha = \frac{\mathbf{x}\cdot\mathbf{y}}{\norm{\mathbf{x}}\,\norm{\mathbf{y}}}
             = \frac{\sum_i x_i y_i}{\sqrt{\sum_i x_i^2}\,\sqrt{\sum_i y_i^2}}.
\]
It is 1 for vectors pointing in the same direction, 0 for orthogonal vectors (with
non-negative counts it is never negative), and it does not change if a vector is
multiplied by a positive number: it ignores frequency.
\end{nbdef}

\begin{nbexample}{dog vs.\ cat and dog vs.\ knife in two dimensions}
\begin{align*}
  \cos(\text{dog}, \text{cat}) &= \frac{115\cdot 52 + 10\cdot 4}{\sqrt{115^2+10^2}\,\sqrt{52^2+4^2}}
     = \frac{6020}{115.43 \cdot 52.15} \approx 1.000 \quad(\alpha \approx 0.6^\circ)\\
  \cos(\text{dog}, \text{knife}) &= \frac{115\cdot 51 + 10\cdot 84}{115.43 \cdot 98.27}
     = \frac{6705}{11343} \approx 0.591 \quad(\alpha \approx 53.8^\circ)
\end{align*}
Now \emph{cat} is (in these two dimensions) practically identical to \emph{dog}, and
\emph{knife} is far. The slides report $\alpha = 54.3^\circ$ for knife--dog; computing it
from the counts gives $53.8^\circ$, the small difference does not matter.
\end{nbexample}

\begin{nbwhy}
Every modern word representation, from count-based DSMs to \textbf{word embeddings}
(\emph{word2vec}, GloVe) to the internal vectors of BERT and GPT, follows this recipe:
a word is a vector learnt from its contexts, and similar words have vectors with a small
angle between them. Word2vec learns short dense vectors (a few hundred dimensions)
directly, instead of counting; implementations in Python and C, with examples, are linked
in the course's AulaWeb module.
\end{nbwhy}

\begin{nbsummary}[distributional semantics]
\begin{itemize}[nosep]
  \item Wittgenstein: meaning = use; Firth: ``you shall know a word by the company it
        keeps''; Harris: difference of meaning correlates with difference of
        distribution.
  \item \emph{bardiwac}: the meaning of an unknown word can be inferred from its
        contexts alone.
  \item Co-occurrence matrix: rows = target words, columns = contexts, cells = counts;
        each row is the word's vector. Weight the counts (e.g.\ $\log(1+f)$) before
        comparing.
  \item Euclidean distance is dominated by frequency; use the \textbf{angle}, i.e.\
        \textbf{cosine similarity} $\mathbf{x}\cdot\mathbf{y}/(\norm{\mathbf{x}}\norm{\mathbf{y}})$,
        or normalise the vectors.
  \item This is the basis of word embeddings (word2vec) and of the representations in
        LLMs.
\end{itemize}
\end{nbsummary}

